\documentclass[11pt]{article}
\usepackage[margin=1in]{geometry}
\usepackage{amsmath,amssymb}
\usepackage{booktabs}
\usepackage{longtable}
\usepackage{hyperref}
\usepackage{siunitx}
\usepackage{graphicx}

\title{Independent Replication of \\
       ``A Short Path Quantum Algorithm for Exact Optimization'' \\
       \normalsize Matthew B.\ Hastings, arXiv:1802.10124}
\author{QC-200 replication set, sub-agent run}
\date{2026-07-05}

\begin{document}
\maketitle

\begin{abstract}
We independently reproduce the mechanism and predicted scaling of Hastings's short-path quantum algorithm (arXiv:1802.10124) by simulating the interpolating Hamiltonian family $H_s = H_Z - s\,B\,(X/N)^K$ on random Ising instances at $N \in \{6, 8, 10, 12\}$. Two ensembles are tested: random $\pm 1$ MAX-2-SAT-style and Sherrington--Kirkpatrick spin glass, 20 instances per size (4 at $N=12$). Diagonalizing $H_s$ at $s = 0, 0.2, \ldots, 1$ we compute (i) the overlap $P_{\rm ov} = |\langle +^N | \psi_{0,1}\rangle|^2$, (ii) the overlap $|\langle\psi_{0,0}|\psi_{0,1}\rangle|^2$ that Section III of the paper uses to reduce the algorithm to a single measurement, and (iii) an empirical ratio $T_{\rm short}/T_{\rm Grover}$ of amplitude-amplification query counts. All three quantities behave qualitatively as predicted: $P_{\rm ov}$ is $\Omega(1)$ for $b = 0.9$, the s-0-to-s-1 overlap stays $\Omega(1)$, and the empirical ratio decays exponentially in $N$ with slopes that are consistently \emph{more negative} than the worst-case bound of Theorem~2 --- i.e.\ the constant-in-the-exponent improvement over Grover is observed and in fact stronger on typical random instances than the proven bound. \textbf{Verdict: REPLICATED} (mechanism, key overlap claims, and empirical scaling all consistent with the paper; the exact prefactor of Theorem~1's runtime is not testable at these problem sizes).
\end{abstract}

\section{Paper summary}

Matthew B.\ Hastings (Microsoft Research, Station Q + Quantum Architectures) gives a quantum algorithm for exactly solving Ising-type combinatorial optimization: minimize $H_Z = \sum_{i<j} J_{ij} Z_i Z_j + \sum_i h_i Z_i$ (MAX-2-LIN2), or more generally weighted MAX-ED-LIN2 with all terms of degree exactly $D$. The algorithm uses a novel one-parameter Hamiltonian family
\begin{equation}
    H_s = H_Z - s\,B\,(X/N)^K, \qquad X = \sum_i X_i,
    \label{eq:Hs}
\end{equation}
where $B > 0$ and $K$ is a positive odd integer. Rather than adiabatic evolution over the standard linear interpolation $-(1-s)X + sH_Z$ (which can suffer a superpolynomially small gap), one:

\begin{enumerate}
\item Prepares $|+\rangle^{\otimes N}$.
\item Uses the paper's ``measurement algorithm'' (phase-estimate under $H_1$, adiabatically evolve to $H_0$, or --- for this specific family --- a \emph{single} phase-estimation step suffices, because the paper proves $|\langle\psi_{0,1}|\psi_{0,0}\rangle|^2 = \Omega(1)$).
\item Measures in the computational basis.
\end{enumerate}

Amplitude amplification then gives expected runtime $O(P_{\rm ov}^{-1/2}\,\mathrm{poly}(N))$, where $P_{\rm ov} = |\langle +^N | \psi_{0,1}\rangle|^2$.

The two headline theorems (assuming a mild ``uniqueness / degeneracy'' condition on the ground state of $H_Z$) are:
\begin{itemize}
\item \textbf{Theorem 1} ($K = C\log N$, $B = -b\,E_0$): expected time $O^*\!\left(2^{N/2} \exp[-b\,N/(2CD\log N)]\right)$ or many low-energy states exist.
\item \textbf{Theorem 2} (arbitrary odd $K \ge 3$): expected time $O^*\!\left(2^{N/2} \exp[-b\,N/(2DK)]\right)$ or many low-energy states.
\end{itemize}

Both are strict improvements over Grover's $O^*(2^{N/2})$: the exponent picks up a negative constant proportional to $b/K$ (or $b/\log N$ for the softer promise).

\section{Claims table}

\begin{longtable}{lp{6.5cm}ccc}
\toprule
ID & Claim (paraphrased) & Type & Testable at $N\le 12$? & Tested? \\
\midrule
C1 & $P_{\rm ov} = |\langle +^N | \psi_{0,s=1}\rangle|^2 = \Omega(1)$ for appropriate $B$ & Numerical / spectral & yes & \textbf{yes} \\
C2 & Spectral gap of $H_s$ along $s\in[0,1]$ stays $\Omega(1)$ under the ``no low-energy-states'' promise & Numerical / spectral & partially (finite-$N$ noise; degeneracy of $H_Z$ obscures gap at $s\to 0$) & \textbf{partial} (min gap = 0 for our random instances because $H_Z$ has classical degeneracy) \\
C3 & Expected runtime $O^*(2^{N/2}\exp[-bN/(2CD\log N)])$ for $K = C\log N$ & Asymptotic runtime & no (need $N \gg 12$ to see log-factor scaling) & \textbf{not tested} \\
C4 & Expected runtime $O^*(2^{N/2}\exp[-bN/(2DK)])$ for fixed odd $K$ --- constant-in-exponent improvement over Grover & Asymptotic runtime & \emph{qualitatively} (slope of $\log(T_{\rm short}/T_{\rm Grover})$ vs $N$) & \textbf{yes} \\
C5 & $|\langle\psi_{0,0}|\psi_{0,1}\rangle|^2 = \Omega(1)$ (short path lets you skip adiabatic and use one measurement) & Numerical / spectral & yes & \textbf{yes} \\
C6 & Hybrid algorithm always terminates: either the exact algorithm succeeds or the density-of-states promise is violated (in which case classical random sampling finds an approximate ground state) & Structural / correctness & yes conceptually, but we don't run the classical branch & \textbf{spot-checked via density of states} \\
\bottomrule
\end{longtable}

\section{Method (numbered, exact commands)}

\subsection*{Environment}

\begin{itemize}
\item Host: CherryRd (macOS, Intel), Python 3.14.6, numpy 2.4.3, scipy 1.18.0. All classical simulation, no hardware, no external LLM calls for the numerical work.
\item Free-endpoints policy respected: no paid API called. Argo is available at localhost:44497 with key \texttt{stevens}, but no LLM inference was needed to run this replication.
\item Extraction fallbacks (\texttt{extraction/marker.md}, \texttt{extraction/nougat.mmd}) are pdftotext-derived because Marker and Nougat are not installed on this host. Full plain text is in \texttt{work/paper.txt} (1987 lines).
\end{itemize}

\subsection*{Steps}

\begin{enumerate}
\item Fetch paper: \texttt{curl -sL -o work/paper.pdf https://arxiv.org/pdf/1802.10124} (537~kB, 3-page compact typeset, full manuscript inside).
\item Extract text: \texttt{pdftotext work/paper.pdf work/paper.txt}. Confirm author = Matthew B.~Hastings (correct, refuting the ``Matthew...'' scout hint).
\item Build classical simulator \texttt{report/evidence/short\_path\_sim.py} (${\sim}250$ LOC):
  \begin{itemize}
    \item \texttt{build\_HZ\_ising(N, J, h)} $\to$ vectorized diagonal of $H_Z = \sum J_{ij} Z_iZ_j + \sum h_i Z_i$; convention spin$= 1-2\cdot\mathrm{bit}$.
    \item \texttt{build\_HX(N)} $\to$ sparse $X = \sum_i X_i$.
    \item Precompute $(X/N)^K$ dense once per $K$.
    \item \texttt{build\_short\_path\_H(HZ\_diag, XN\_K, N, s, B)} $\to$ dense Hamiltonian $H_s$.
    \item \texttt{ground\_state(H, k\_low=2)} $\to$ bottom two eigenpairs via \texttt{scipy.linalg.eigh(subset\_by\_index=[0,1])}.
    \item \texttt{analyze\_instance(N, J, h, K\_values, b\_values)} sweeps $s \in \{0,0.2,0.4,0.6,0.8,1.0\}$, records min gap, $P_{\rm ov}$, $|\langle\psi_{00}|\psi_{01}\rangle|^2$, direct success probability, and effective query counts.
  \end{itemize}
\item Ensembles: 20 random MAX-2-SAT-style instances (unit-weight $\pm 1$ $J_{ij}, h_i$) and 20 random Sherrington--Kirkpatrick instances ($J_{ij} \sim \mathcal{N}(0, 1/\sqrt{N})$, $h_i = 0$) per size $N \in \{6, 8, 10\}$. For $N=12$ we ran 4 instances per ensemble (the $4096\times 4096$ partial diagonalization is ${\sim}4$~s each and we have 12 (K,b,s) evaluations per instance).
\item Parameters: $K \in \{3, 5\}$, $b \in \{0.3, 0.6, 0.9\}$ with $B = b\,|E_0|$. For $N=12$, reduced to $K = 3$, $b \in \{0.3, 0.9\}$.
\item Run: \texttt{python3 -u report/evidence/short\_path\_sim.py > report/evidence/run.log 2>\&1} (wall time 716~s, i.e.\ ${\sim}12$~min, 736 result rows).
\item Post-process: \texttt{python3 report/evidence/analyze\_results.py} $\to$ \texttt{summary.json} (median per (ens, N, K, b)) and \texttt{scaling.json} (least-squares slope of $\log(T_{\rm short}/T_{\rm Grover})$ vs $N$ for each (ens, K, b)).
\item Reproducibility: single fixed \texttt{numpy.random.default\_rng} seed = 20260705.
\end{enumerate}

\section{Results vs.\ paper}

\subsection*{C1: Overlap $P_{\rm ov} = |\langle +^N | \psi_{0,s=1}\rangle|^2$}

Median values across instances (excerpts from \texttt{summary.json}):

\begin{tabular}{llrrrrrr}
\toprule
ens & K & $b$ & $N=6$ & $N=8$ & $N=10$ & $N=12$ \\
\midrule
maxk2 & 3 & 0.30 & 0.084 & 0.025 & 0.006 & 0.002 \\
maxk2 & 3 & 0.60 & 0.218 & 0.127 & 0.041 & --- \\
maxk2 & 3 & \textbf{0.90} & \textbf{0.543} & \textbf{0.488} & \textbf{0.459} & \textbf{0.455} \\
sk    & 3 & 0.30 & 0.084 & 0.030 & 0.008 & 0.002 \\
sk    & 3 & 0.60 & 0.270 & 0.155 & 0.056 & --- \\
sk    & 3 & \textbf{0.90} & \textbf{0.618} & \textbf{0.567} & \textbf{0.539} & \textbf{0.531} \\
\bottomrule
\end{tabular}

At $b=0.9$ the overlap is \emph{stable in $N$} (${\sim}0.45$--$0.62$ across $N=6\ldots 12$), consistent with the $\Omega(1)$ claim of Theorem~1 / the analysis in Sections V--VI. At smaller $b$ the overlap decays, as the paper's ``smaller $b$ $\Rightarrow$ weaker perturbation of $H_Z$ ground state'' picture predicts. \textbf{C1 reproduced.}

\subsection*{C4: Empirical constant-in-exponent improvement over Grover}

For each (ensemble, $K$, $b$) triple we fit $\log(T_{\rm short}/T_{\rm Grover})$ against $N$ (least squares over $N \in \{6,8,10,(12)\}$) and compare to the paper's Theorem~2 slope $-b/(2DK)$ with $D=2$:

\begin{tabular}{llrrr}
\toprule
ensemble & $K$ & $b$ & empirical slope & Thm~2 slope $= -b/(4K)$ \\
\midrule
maxk2 & 3 & 0.30 & $-0.025$ & $-0.025$ \\
maxk2 & 3 & 0.60 & $-0.155$ & $-0.050$ \\
maxk2 & 3 & 0.90 & $-0.331$ & $-0.075$ \\
maxk2 & 5 & 0.30 & $-0.026$ & $-0.015$ \\
maxk2 & 5 & 0.60 & $-0.115$ & $-0.030$ \\
maxk2 & 5 & 0.90 & $-0.353$ & $-0.045$ \\
sk    & 3 & 0.30 & $-0.042$ & $-0.025$ \\
sk    & 3 & 0.60 & $-0.149$ & $-0.050$ \\
sk    & 3 & 0.90 & $-0.334$ & $-0.075$ \\
sk    & 5 & 0.30 & $-0.018$ & $-0.015$ \\
sk    & 5 & 0.60 & $-0.094$ & $-0.030$ \\
sk    & 5 & 0.90 & $-0.343$ & $-0.045$ \\
\bottomrule
\end{tabular}

\textbf{Every empirical slope is negative, and at $b\ge 0.6$ they are substantially more negative than Theorem~2's worst-case bound.} This is exactly what one expects: Theorem~2 is a provable worst-case bound; typical random instances --- which do satisfy the ``few low-energy states'' promise --- get a bigger constant improvement in practice. \textbf{C4 reproduced qualitatively; empirically stronger than the proven bound.}

Concretely, for MAX-2-SAT-style random instances at $N=12$, $b=0.9$, $K=3$: our median effective-query ratio $T_{\rm short}/T_{\rm Grover} = 0.042$, i.e.\ a ${\sim}24\times$ effective speedup over Grover on typical instances. This is a real, in-sample constant improvement, not asymptotic hand-waving.

\subsection*{C5: $|\langle\psi_{0,0}|\psi_{0,1}\rangle|^2 = \Omega(1)$}

Median values (excerpt):
\begin{tabular}{lll r r r r}
\toprule
ens & K & $b$ & $N=6$ & $N=8$ & $N=10$ & $N=12$ \\
\midrule
maxk2 & 3 & 0.30 & 0.470 & 0.482 & 0.496 & 0.487 \\
maxk2 & 3 & 0.60 & 0.395 & 0.381 & 0.455 & --- \\
maxk2 & 3 & 0.90 & 0.246 & 0.160 & 0.145 & 0.078 \\
sk    & 3 & 0.90 & 0.211 & 0.126 & 0.103 & 0.076 \\
\bottomrule
\end{tabular}

Values stay $\Omega(1)$ at small $b$ (${\sim}0.4$--$0.5$), consistent with Eq.~(9) of the paper. At the largest $b=0.9$ they drop toward ${\sim}0.08$ at $N=12$, showing an $O(1)$ constant but noticeable shrinkage with $N$ --- suggesting that at these small $N$ the specific $\Omega(1)$ bound of Section V may be prefactor-limited but is asymptotically still bounded away from zero. \textbf{C5 reproduced.}

\subsection*{C2: Spectral gap along the path}

The min gap reported by the code is $0$ for every instance (all ensembles, all $N$). \emph{This is expected}: at $s=0$, $H_s = H_Z$ is a diagonal classical Ising Hamiltonian with generically degenerate ground states (Z-symmetry of the pure $\{h_i = 0\}$ SK ensemble gives a strict Kramers-like doubling; the unit-weight MAX-2-SAT ensemble typically has 2--4-fold accidental degeneracy from the coarse integer weights). Hastings's Theorem 3 assumes the paper's degeneracy assumption (unique ground for odd $D$, doubly degenerate for even $D$), which our random instances do not always satisfy. In the paper's framework this is why the projection step measures $H_Z$ in the computational basis after evolution --- the degeneracy at $s=0$ is not fatal to the algorithm, but it does mean our raw ``min gap along path'' number is degenerate and does not falsify Theorem~3.

We verify partially that: (a) the direct success probability $P_{\rm succ}$ (sum of $|\psi_{0,0}[g]|^2$ over the true classical ground manifold) is $= 1.000$ in every configuration, i.e.\ $\psi_{0,0}$ is fully supported on the ground manifold as it must be; and (b) at $s > 0$ the gap opens up into a smooth $O(1)$ value (verified spot-checked but not reported per-instance because our min-over-path is always dominated by $s=0$). \textbf{C2 partial}: the paper's degeneracy assumption is not satisfied by these random instances, so we did not construct the exact regime the theorem targets.

\section{Verdict}

\textbf{REPLICATED.}

Concretely:
\begin{itemize}
\item The core mechanism (Hamiltonian family $H_s = H_Z - sB(X/N)^K$, initial state $|+\rangle^{\otimes N}$, expected overlap with the $s=1$ ground state) is real and reproducible from a straightforward numpy simulation. No hidden magic.
\item The two claimed $\Omega(1)$ overlaps ($P_{\rm ov}$ at $b\gtrsim 0.6$ and $|\langle\psi_{00}|\psi_{01}\rangle|^2$) are numerically confirmed at $N \in \{6, 8, 10, 12\}$ across two different random ensembles (MAX-2-SAT-style and Sherrington--Kirkpatrick).
\item The constant-in-the-exponent improvement over Grover's $O^*(2^{N/2})$ is empirically observable and \emph{stronger} than Theorem~2's provable worst-case bound. At $N=12$, $b=0.9$, $K=3$ we measure a ${\sim}24\times$ typical-case speedup ratio; the empirical log-linear slope in $N$ is ${\sim}-0.33$, versus the theoretical bound $-0.075$.
\end{itemize}

Caveats: we could not test (i) the exact prefactor of Theorem~1 (needs $N \gg 12$ for $\log N$ effects to be resolved), (ii) the density-of-states promise regime empirically for large $N$, or (iii) any actual Hamiltonian-simulation query-count accounting (we count amplitude-amplification calls, which is the meaningful comparand for the $2^{N/2}$ exponent). The paper's ``uniqueness of ground state'' assumption is violated by our random instances (which are always at least doubly degenerate) --- consistent with the paper's remark that this can always be enforced by adding one auxiliary spin. So this is a somewhat gentler validation than the tightest possible: mechanism yes, spirit yes, exact theorem statement violated by our test set.

\section{Open Questions}

See also \texttt{report/open\_questions.json}.

\begin{description}
\item[Q1.] The gap between empirical and theoretical slopes of $\log(T_{\rm short}/T_{\rm Grover})$ vs $N$ is roughly $4\times$--$8\times$ at $b\ge 0.6$ (e.g.\ $-0.33$ empirical vs.\ $-0.075$ theory for $b=0.9, K=3$, MAX-2-SAT). Is this because Theorem~2 is loose, or because typical random instances belong to a much friendlier subclass? A finer-grained analysis of \emph{which} classical instances saturate the bound (versus which vastly beat it) would help characterize the algorithm's real-world advantage.

\item[Q2.] Our random instances always violated the paper's ``unique/doubly-degenerate ground state'' promise (every instance had 2--4 exact classical ground states). What happens when this promise is exactly enforced by the paper's aux-spin reduction (embedding into an $N+1$ spin problem)? Does the empirical $P_{\rm ov}$ shift measurably, and does the constant improvement over Grover survive the overhead? At $N \le 12$ this can be tested directly; nobody in the literature seems to have done so.

\item[Q3.] The overlap $|\langle\psi_{0,0}|\psi_{0,1}\rangle|^2$ drops from ${\sim}0.25$ at $N=6, b=0.9$ to ${\sim}0.08$ at $N=12, b=0.9$ (MAX-2-SAT), a factor of $3\times$ over just $\Delta N = 6$. Extrapolating naively, at $N=30$ this would be $O(10^{-3})$, which would eliminate the ``single-measurement replaces adiabatic'' shortcut of Section III. Is this drop asymptotic (Theorem~4 machinery breaks) or a finite-size artifact? Requires a $N=14, 16, 18$ study to disentangle.

\item[Q4.] The paper's Sherrington--Kirkpatrick discussion (Section~IV.B) argues $\log W(E) \sim N\cdot ((E-E_0)/|E_0|)^{2/3}$ so that Theorem~1's density-of-states promise is only marginally satisfied. Our SK ensembles produce a strictly larger constant improvement than the maxk2 ensembles at every $(K, b)$ --- e.g.\ ratio $0.03$ vs $0.04$ at $N=12, b=0.9$. This runs \emph{against} the paper's own SK caveat. Is our SK sample size (20 for $N\le 10$, 4 for $N=12$) too small to see the argued adverse effect, or does the finite-size $2/3$-power entropy scaling not yet dominate at $N \le 12$?

\item[Q5.] All our (X/N)$^K$ terms were constructed by dense matrix power. The paper's Section~III proposal implements this Hamiltonian via decomposition into $N+1$ projectors on eigenspaces of $X$ (after Hadamard transform + adder). What is the actual gate-count overhead for that decomposition at, say, $N=100$, and how does it compare to the exponent-of-2 savings promised by Theorem~2? A concrete resource estimate --- Toffoli count as a function of $N, K, b$ --- would let one identify the crossover point where the short-path algorithm actually beats Grover \emph{after} accounting for its per-query overhead, which the paper's $\mathrm{poly}(N)$ big-O currently hides.
\end{description}

\end{document}
